% --- 1. INFORMAZIONI GENERALI ---
\section{Informazioni Generali}
\begin{itemize}
	\item \textbf{Azienda Coinvolta:} {Zucchetti}
	\item \textbf{Mezzo di Comunicazione:} {Google Meet}
	\item \textbf{Data:} 2026-03-25
	\item \textbf{Orario di inizio -- fine:} 17:00 -- 17:25
\end{itemize}


\vspace{0.5cm}
\textbf{Referenti Azienda:}
\begin{table}[ht]
	\centering
	\renewcommand{\arraystretch}{1.2}
	\begin{tabularx}{\textwidth}{X l}
		\toprule
		\textbf{Nome e Cognome} & \textbf{Ruolo} \\
		\midrule
		{Gregorio Piccoli} & {Referente Principale} \\
		\bottomrule
	\end{tabularx}
\end{table}

\vspace{0.5cm}
\textbf{Partecipanti Gruppo (Synergo Unit):}
\begin{table}[ht]
	\centering
	\renewcommand{\arraystretch}{1.2}
	\begin{tabularx}{\textwidth}{X c}
		\toprule
		\textbf{Nome e Cognome} & \textbf{Matricola}\\
		\midrule
		Sofia De Blasi & 2111014 \\
		Roberto Mariano Doroftei & 2111031 \\
		Filippo Idiometri & 2035617 \\
		Francesco Marcon & 2110990 \\
		Armando Moda Scarati & 2082864 \\
		Anton Ricardo Rupi & 2054304 \\
		\bottomrule
	\end{tabularx}
\end{table}

% --- 2. ORDINE DEL GIORNO ---
Gli argomenti previsti per la riunione sono:
\begin{enumerate}[leftmargin=1.8em, itemsep=3pt]
    \item Formazione e supporto tecnico da parte dell'azienda
    \item Disponibilità di esempi o test esistenti per i ``Sei Cappelli per Pensare\textsubscript{G}''
    \item Separazione tra funzionalità base e applicazione dei Sei Cappelli
    \item Struttura e obiettivo dell'output dei Sei Cappelli
    \item Requisiti di interfaccia e usabilità\textsubscript{G}
    \item Chatbot\textsubscript{G} con messaggi precompilati\textsubscript{G} e accessibilità
\end{enumerate}

% --- 3. RESOCONTO ---
\section{Resoconto}

\subsection{Formazione e Supporto Tecnico da Parte dell'Azienda}

\textbf{Discussione:}\\
Il gruppo ha chiesto se siano previsti incontri strutturati di formazione su tematiche come le metodologie di project design: definizione degli obiettivi, scelta degli strumenti e delle tecnologie.

\vspace{0.3cm}
\textbf{Risposta e accordi con il proponente:}
\begin{itemize}[leftmargin=1.5em, itemsep=3pt]
    \item Il proponente ha confermato disponibilità a fornire supporto \textit{su richiesta} laddove necessario, in particolare riguardo all'invocazione di LLM\textsubscript{G} (chiamate API\textsubscript{G}, struttura dei messaggi, parametri principali).
    \item Per le tecnologie da adottare, l'azienda concede \textbf{piena libertà} al gruppo: è possibile esplorare strumenti già noti o nuovi, senza vincoli imposti.
\end{itemize}

\vspace{0.1cm}
\textbf{Impegni del gruppo:}
\begin{itemize}[leftmargin=1.5em, itemsep=3pt]
    \item Segnalare al proponente eventuali necessità formative.
\end{itemize}

\vspace{0.4cm}

\subsection{Disponibilità di Esempi o Test Esistenti — ``Sei Cappelli per Pensare''}

\textbf{Discussione:}\\
Il gruppo ha chiesto se il proponente disponga di esempi o casi di test già esistenti da utilizzare come riferimento per migliorare e valutare i prompt\textsubscript{G} relativi all'implementazione dei ``Sei Cappelli per Pensare''.

\vspace{0.3cm}
\textbf{Risposta e accordi con il proponente:}
\begin{itemize}[leftmargin=1.5em, itemsep=3pt]
    \item Non sono disponibili esempi o test predefiniti.
    \item L'approccio atteso è quello dell'esplorazione sperimentale del \textbf{prompt engineering\textsubscript{G}}, verificando autonomamente fino a che punto il modello linguistico sia in grado di applicare la metodologia dei cappelli in modo efficace.
    \item La valutazione è qualitativa e iterativa, basata sull'osservazione diretta dei risultati prodotti durante i test.
\end{itemize}

\vspace{0.4cm}

\subsection{Separazione tra Funzionalità Base e Applicazione dei "Sei Cappelli"}

\textbf{Discussione:}\\
Il gruppo ha chiesto se i "Sei Cappelli" debbano essere applicati internamente alle funzionalità base (per esempio, eseguire un riassunto già filtrato attraverso il cappello nero (critico)) oppure se debbano costituire un insieme separato di prompt.

\vspace{0.3cm}
\textbf{Risposta e accordi con il proponente:}
\begin{itemize}[leftmargin=1.5em, itemsep=3pt]
    \item Il proponente ha chiarito che i due ambiti devono essere \textbf{distinti}:
    \begin{enumerate}[leftmargin=1.8em, itemsep=2pt]
        \item \textit{Funzionalità base}: generazione di testo, riassunto, traduzione, riscrittura di porzioni di testo gestite con prompt propri.
        \item \textit{Sei Cappelli}: applicabili \textbf{esclusivamente alla generazione autonoma} di una parte del testo, tramite prompt separati e dedicati.
    \end{enumerate}
    \item Combinare i cappelli con ciascuna funzionalità base richiederebbe una matrice $5 \times 6$ di prompt, rendendo il sistema molto più complesso.
    \item Il flusso corretto è \textbf{sequenziale}: prima si applica una funzionalità base (es.\ riassunto), poi, separatamente e su richiesta, si applicano i cappelli.
\end{itemize}

\vspace{0.1cm}
\textbf{Impegni del gruppo:}
\begin{itemize}[leftmargin=1.5em, itemsep=3pt]
    \item Progettare l'architettura dei prompt in due blocchi indipendenti, mantenendo separati i casi d'uso base dai Sei Cappelli.
\end{itemize}

\vspace{0.4cm}

\subsection{Struttura e Obiettivo dell'Output dei Sei Cappelli}

\textbf{Discussione:}\\
Il gruppo ha chiesto se uno degli output attesi fosse quello di generazione di un report finale in PDF contenente le sei analisi corrispondenti ai cappelli, chiedendo conferma sull'approccio e sullo sviluppo dei prompt correlati.

\vspace{0.3cm}
\textbf{Risposta e accordi con il proponente:}
\begin{itemize}[leftmargin=1.5em, itemsep=3pt]
    \item Il proponente ha \textbf{sconsigliato} la generazione di un report PDF: l'obiettivo è fornire un'analisi \textit{contestuale e immediata} del testo in fase di scrittura, non un documento separato.
    \item Il meccanismo preferito è: \textbf{selezione del testo} → bottone corrispondente al cappello desiderato → risposta visualizzata in un \textit{pop-up\textsubscript{G}} o pannello laterale.
    \item In alternativa, è valutabile un approccio \textit{chatbot} che analizza su richiesta l'intero testo o una parte selezionata, ma rimane un requisito opzionale.
    \item L'idea base rimane: \textbf{un bottone per cappello}, con risposta immediata e non persistente sotto forma di report.
    \item Durante la riunione il proponente ha mostrato una presentazione che distingue chiaramente due concetti:
    \begin{itemize}[leftmargin=1.5em, itemsep=2pt]
        \item \textit{Bottone/azione}: esegue una funzionalità specifica e ben definita.
        \item \textit{Chatbot}: l'utente esprime liberamente un'intenzione/un \textbf{compito\textsubscript{G}}  (es.\ ``voglio applicare il cappello critico'') e il chatbot individua la funzionalità corrispondente, agendo da \textbf{mediatore tra linguaggio naturale e funzionalità}.
    \end{itemize}
\end{itemize}

\vspace{0.1cm}
\textbf{Impegni del gruppo:}
\begin{itemize}[leftmargin=1.5em, itemsep=3pt]
    \item Valutare il chatbot come requisito opzionale aggiuntivo, senza sostituire l'interfaccia base a bottoni.
\end{itemize}

\vspace{0.4cm}

\subsection{Requisiti di Interfaccia e Usabilità}

\textbf{Discussione:}\\
Il gruppo ha chiesto conferma che la semplicità e l'intuitività dell'interfaccia siano i requisiti prioritari, e come il chatbot si inserisca in questo contesto.

\vspace{0.3cm}
\textbf{Risposta e accordi con il proponente:}
\begin{itemize}[leftmargin=1.5em, itemsep=3pt]
    \item Il proponente ha confermato che \textbf{semplicità e immediatezza} sono requisiti fondamentali dell'interfaccia.
    \item Il chatbot, se implementato, deve essere un elemento \textit{complementare}: non deve sovrapporsi o complicare l'interfaccia base.
    \item L'interfaccia principale deve rimanere funzionale e comprensibile anche senza il chatbot.
\end{itemize}

\vspace{0.4cm}

\subsection{Chatbot con Messaggi Precompilati e Accessibilità}

\textbf{Discussione:}\\
Il gruppo ha proposto di rendere il chatbot più deterministico tramite l'uso di \textbf{messaggi precompilati}: una volta che il chatbot individua l'intenzione dell'utente, propone bottoni con azioni già scelte (es.\ ``Vuoi applicare la critica?'', ``Vuoi un riassunto?''), riducendo l'ambiguità e guidando la conversazione.

\vspace{0.3cm}
\textbf{Risposta e accordi con il proponente:}
\begin{itemize}[leftmargin=1.5em, itemsep=3pt]
    \item Il proponente ha \textbf{approvato} l'approccio: i messaggi precompilati rendono il flusso più prevedibile e controllabile.
    \item È stato sottolineato il valore aggiunto in termini di \textbf{accessibilità}: bottoni contestuali all'interno della chat riducono il carico cognitivo e facilitano l'uso da parte di utenti con difficoltà nell'interazione testuale libera.
\end{itemize}

\vspace{0.1cm}
\textbf{Impegni del gruppo:}
\begin{itemize}[leftmargin=1.5em, itemsep=3pt]
    \item Includere i messaggi precompilati nell'eventuale progettazione del chatbot come pattern di interazione principale.
\end{itemize}

% --- 4. DECISIONI FORMALI ---
\section{Decisioni Formali}
Nella seguente tabella sono tracciate le decisioni ufficiali e le azioni conseguenti approvate congiuntamente all'azienda.

\renewcommand{\arraystretch}{1.4}

\begin{longtable}{c p{10.5cm} c l}
\toprule
\textbf{Cod.} & \textbf{Decisione e Azione} & \textbf{Impatto} & \textbf{Scadenza} \\
\midrule
\endfirsthead

\toprule
\textbf{Cod.} & \textbf{Decisione e Azione} & \textbf{Impatto} & \textbf{Scadenza} \\
\midrule
\endhead

\midrule
\multicolumn{4}{r}{\textit{Continua nella pagina successiva}} \\
\midrule
\endfoot

\bottomrule
\endlastfoot

\textbf{VE-3.1} &
\textbf{Decisione:} Le funzionalità base (riassunto, traduzione, riscrittura) e i Sei Cappelli sono implementati come insiemi di prompt separati e indipendenti.\newline
\textbf{Azione:} Progettare l'architettura dei prompt in due blocchi distinti. &
Requisiti & Da concordare \\

\textbf{VE-3.2} &
\textbf{Decisione:} L'output dei Sei Cappelli sarà visualizzato in pop-up contestuale, non in report PDF.\newline
\textbf{Azione:} Implementare meccanismo di selezione testo + bottone per cappello. &
Requisiti & Da concordare \\

\textbf{VE-3.3} &
\textbf{Decisione:} Il chatbot è un requisito opzionale; se implementato, utilizza messaggi precompilati per guidare l'utente.\newline
\textbf{Azione:} Inserire il chatbot con bottoni contestuali come feature opzionale nel backlog. &
Requisiti & Da concordare \\

\end{longtable}

% --- 5. PROSSIMO INCONTRO ---
\section{Prossimo Incontro}
\begin{itemize}
	\item \textbf{Data e Ora:} {Da concordare}
	\item \textbf{Modalità/Luogo:} {Google Meet}
	\item \textbf{Argomenti Previsti (in caso di aggiudicazione):}
	\begin{itemize}
		\item Presentazione della proposta architetturale e dello stack tecnologico\textsubscript{G} scelto
        \item Revisione\textsubscript{G} del backlog iniziale e delle priorità di sviluppo del PoC\textsubscript{G}
        \item Eventuali chiarimenti emersi dalla progettazione dei prompt
	\end{itemize}
\end{itemize}

\newpage
% --- 6. FIRME ---
\section*{Approvazione e Firme}
Con le seguenti firme, i rappresentanti approvano formalmente il contenuto del presente verbale, i requisiti concordati e le decisioni prese.

\vspace{2.5cm}

\noindent
\begin{minipage}[c]{0.45\textwidth}
	\begin{center}
		\rule{6cm}{0.4pt} \\ 
		\vspace{0.2cm}
		\textbf{Per il Gruppo Synergo Unit} \\
		{[Sofia De Blasi]} \\
		\textit{Responsabile di Progetto}
	\end{center}
\end{minipage}
\hfill
\begin{minipage}[c]{0.45\textwidth}
	\begin{center}
		\rule{6cm}{0.4pt} \\ 
		\vspace{0.2cm}
		\textbf{Per l'Azienda Zucchetti} \\
		{[Gregorio Piccoli]} \\
		\textit{Referente Aziendale}
	\end{center}
\end{minipage}

\vspace{2cm}
